\section{Restriction}\label{sec:restrict}

A warp or a thread uses one slice of an address map: the map with its own
warp and thread coordinates fixed.

\begin{definition}[Restriction]\label{def:restrict}
Let $f : \mathrm{Coord}(S) \to \N$ with $S = (n_1,\dots,n_m)$ flat, let
$F \subseteq \{1,\dots,m\}$, and let $p_i \in [0,n_i)$ for each $i \in F$.
The \emph{restriction} of $f$ to $p$ is the map $f|_p$ on the coordinates of
the shape $(n_i)_{i \notin F}$ given by
\[
f|_p(c) = f(c'), \qquad
c'_i = \begin{cases} p_i & i \in F,\\ c_i & i \notin F. \end{cases}
\]
\end{definition}

$f|_p$ is in general not a layout: $f|_p(\mathbf{0})$ is the contribution
of the fixed entries, which need not be $0$, while every layout vanishes at
$\mathbf{0}$.

$f|_p$ is in general not dense, since its image omits the other slices, so
it is not a valid map into $\logical$ or $\tv$ (Definition~\ref{def:validity})
and cannot be the right operand of a composition
(Definition~\ref{def:compose}). Restriction is therefore allowed only for
maps into $\physical$, and is applied after composition, where the warp and
thread coordinates are entries of the composite's domain. Restricting a
restriction is restricting by the union of the fixed entries.

The address expression of a slice is Section~\ref{sec:decide} applied to
$f|_p - f|_p(\mathbf{0})$, which vanishes at the origin, plus the constant
$f|_p(\mathbf{0})$; when that map is not a layout, the composite's expression
is emitted with the fixed coordinates replaced by their values. The difference
can take negative values, which is why Section~\ref{sec:decide} is stated
over $\Z$: for $(2,2){:}(2,1)$ under the swizzle with $b = 1$, $\mathit{src} = 1$, $\mathit{dst} = 0$
(Definition~\ref{def:swizzle}), that is $\sigma(x) = x \oplus ((x \gg 1) \wedge 1)$,
$f(a,b) = (2a+b) \oplus a$, fixing $a = 1$ gives $f|_p(b) = 3 - b$ and
$f|_p - f|_p(\mathbf{0}) = -b$. The emitted arithmetic then carries a negative
coefficient. For
$f = (4,8,4){:}(32,1,8)$, $F = \{1\}$ and $p_1 = 2$, the slice is
$f|_p(r,v) = 64 + r + 8v$.
